\section{在线算法：请求路由与 KV 交接}

\subsection{从路径选择到实例对选择}
PDblend 为请求维护输入长度 $n$、输出预算 $m$、已输出 token 数以及首 token 和完成时刻。
在线调度包含两种路径：混合路径 M 在一个实例内完成预填充与解码；分离路径 PD 选择有序实例对 $(P,D)$，分别执行两个阶段。
池配置发布各实例角色与长度阈值后，路由器按照该阈值生成当前路径的候选，长输入倾向 PD，短输入倾向 M；没有可用 M 时也可使用 PD。
压力变化通过后续配置提交更新路由阈值，已经接纳的请求保持原实例绑定。

PD 候选必须是两个不同实例，且模型、TP、PP、池标识与拓扑代次一致，当前要求 PP 为 1。
这使接收端可以直接使用相同形状的 KV 分片，避免把跨 TP 重排隐含为零成本。
基础负载规则优先选择待预填充 token 较少的 P，再选择在途序列较少的 D；M 优先选择在途序列较少的实例。
启用 SLO 路由时，系统枚举当前路径的全部候选，并为 PD 请求附加 M 备选。
对每个候选涉及的实例，已拥有请求与新请求的输入和最大输出之和不得超过已测得的 KV 容量，同时检查上下文长度和接单状态。
这里的预算是代理侧的保守准入记账，由引擎完成实际 KV block 分配。

随后，路由器根据当前频率、P 上尚未产生首 token 的请求、D 上在途序列数以及上下文长度查询 profiler。
等待成本按各请求的预填充成本求和，避免将多个 prompt 拼成超出测量范围的虚构大请求。
设 $q_P$ 为上述等待时间，$P(n,f_P)$ 为新请求的预填充时间，$D(b,c,f_D)$ 为解码步长。
普通 M 路径以等待、预填充和一个解码步预测首 token；启用 deadline safety 时，还通过虚拟运行槽位估计 FIFO 等待，并计入新预填充对已有解码的暂停。
这些估计来自代理观测，尚不等同于原生调度器队列遥测。

SLO 筛选要求新请求的 TTFT、TPOT 不超过安全系数乘以相应目标，默认安全系数为 $0.85$。
向 M 分流时还检查已有请求的首 token 截止时间和剩余输出预算，防止新预填充损害原有解码。
先在原路径的安全候选中选预测 TTFT 最小者；仅当 PD 没有安全候选时，才选择安全的 M。
若没有可证明安全的选择，当前实现允许沿容量可行的原路径继续服务，并明确记录 SLO 未被证明；容量也不足时拒绝接纳。
因此，这是受配置约束、带保护性分流的路由，而非每次到达都在所有路径间求全局最优。

\subsection{首 token 协议与传输成本}
PD 路径首先向 P 请求一个输出 token，并取得其权威 token ID。
代理立即将该 token 返回客户端，再将原 prompt 与该 token 一并提交给 D，令 D 生成剩余 $m-1$ 个 token。
P2P connector 根据包含两端地址的请求标识传递 KV，代理不搬运 KV tensor；D 复用原 prompt 的 KV，并在本地处理追加的 token。
这一过程避免丢弃 P 已生成的首 token，也保持输出总预算不变。
当 $m=1$ 时，直接由 P 完成请求，无须远程 KV。
当前 carry 协议适用于 token-ID 输入、贪心解码及固定输出预算的受控执行路径。

KV 开销应按照用户看到的输出时序计费。
令 $h$ 表示 P 首 token 到 D 首 token 的间隔，则端点感知路由使用
\begin{align}
 T_{\mathrm{first}}^{PD}&=q_P+P(n,f_P),\\
 T_{\mathrm{token}}^{PD}&=\frac{h+(m-2)D(b,c,f_D)}{m-1},\quad m\ge2.
\end{align}
首间隔覆盖尚未完成的 KV 传递、协议处理、接收端加载及 D 首步；由于 connector 可以在 P 执行中按层提交数据，不应把全部物理复制简单置于首 token 之后。
若 profiler 提供该输入、输出、频率和负载域的端点测量，则直接使用其首间隔预测，避免重复加上已经包含的 D 首步。
较长输出缺少端点覆盖时，现有近似采用
\begin{equation}
 T_{\mathrm{copy}}=t_0+\frac{n b_{KV}}{B},\qquad
 h=\max\{T_{\mathrm{copy}}+D(b,c,f_D),h_{\min}\},
\end{equation}
其中 $b_{KV}$ 是每 token 的 KV 字节数，$B$ 是有效带宽，$t_0$ 为固定成本。
两 token 请求只有一个输出间隔，无法摊薄交接延迟，因此缺少对应端点覆盖时不能被标记为满足 SLO；若有安全 M，则转到 M。

请求执行期间，系统记录实际首间隔、TTFT 和输出进度，供风险监控与后续调度使用。
正常完成后释放记账；失败或客户端断开时，若原生执行状态仍不确定，则保留请求所有权并暂停相关实例接单，收到包含拓扑代次和各 rank 清理证据的取消确认后才释放资源。

\paragraph{跨池评分的适用范围。}
常驻不同 TP 池时，外层按目标流量份额选择可行池，再由池内路径执行。
可选的逐请求能量排序要求所有候选均具有相同测量窗口下的增量能量覆盖；该能量由带请求与基线的全路由 GPU 能量之差得到，传输成本已进入整体测量。
覆盖不全时，整个候选集合回退到时延评分，不能把一般功率拟合视为已验证的请求增量能量。
当前端点首间隔模型主要用于普通 Router 的 SLO 路由；池级规划和 resident 评分仍保留将传输加到首 token 估计的旧近似，尚未形成统一的端点预算模型。

% 可选算法，主稿可改为 algorithmic 或紧凑的 numbered box。
\begin{quote}\small
\textbf{算法：带交接预算的请求调度}\par
1. 根据当前阈值枚举 M 实例或兼容 $(P,D)$ 对；PD 另列 M 备选。\par
2. 过滤长度、KV 预算和接单状态；预测候选时延并检查已有请求预算。\par
3. 原路径有安全候选则选最小 TTFT；否则 PD 请求选择安全 M。\par
4. 若仅原路径容量可行，保留该路径并标记 SLO 未证实；否则拒绝。\par
5. 固定实例与代次并记账；M 直接执行，PD 执行 P 首 token 与 D 续写。\par
6. 完成后释放；清理未确认的失败保留所有权，等待原生回执。
\end{quote}
